\section{Límites de la evidencia, Q\&A y sistemas de investigación de ciclo cerrado}

Después del minuto 58, la imagen muestra principalmente al ponente y al público, pero el contenido no es una charla informal que pueda omitirse. Las preguntas y respuestas cuestionan la escala de los datos, los controles sanos, la validación causal, la conversión comercial y el AI scientist. Esta parte vuelve a situar los atractivos diagramas de modelos anteriores frente a las restricciones reales, y es decisiva para juzgar si la nota es confiable.

\subsection{Tamaño de muestra, controles sanos y problemas fuera de distribución}

El público pregunta primero si el número de pacientes es suficiente. La respuesta no se centra en dar un umbral universal, sino en distinguir tareas: la reconstrucción a nivel celular puede beneficiarse de una gran cantidad de mediciones locales, mientras que la extrapolación de la eficacia a nivel de paciente sigue limitada por el número de pacientes independientes y de etiquetas de intervención. Otra carencia fundamental es que las personas sanas normalmente no se someten a biopsias tumorales, por lo que la distribución de entrenamiento está sesgada de origen hacia tejido enfermo.

\teachervoice{00:58:43--01:02:52, el docente reconoce que tanto el tamaño de muestra patient-level como los controles sanos son limitaciones. La falta de tejido sano no solo afecta las preguntas de prevención, sino que también dificulta que el modelo aprenda la transición completa de estados desde el tejido sano hasta el canceroso.}

\begin{warningbox}{Missing not at random}
La ausencia de datos sanos no es una ausencia aleatoria, sino que la determina el flujo clínico: a las personas sanas normalmente no se les toma tejido. Esto significa que una imputación simple no puede recuperar la distribución de la población no observada; se necesitan cohorts externas, mediciones poco invasivas, datos longitudinales o un modelo explícito del sesgo de selección.
\end{warningbox}

\subsection{De la predicción del modelo a la causalidad biológica y el valor clínico}

La cadena de descubrimiento de fármacos incluye al menos target hypothesis, intervención experimental, revisión del mecanismo, toxicidad, seguridad y eficacia clínica. Un world model puede reducir el espacio de candidatos, pero no puede saltarse las etapas posteriores. Si el modelo sugiere una diana, el siguiente paso más valioso normalmente no es generar más texto explicativo, sino diseñar experimentos que distingan entre mecanismos competidores.

\teachervoice{01:03:21--01:06:01, la discusión en clase separa «descubrir una diana» de «demostrar que es segura y eficaz». El docente enfatiza que la salida del modelo sigue requiriendo experimentos reales, y que el éxito comercial o clínico no puede deducirse directamente de predicciones virtuales.}

\begin{importantbox}{La validación más sólida es una predicción que puede fallar}
Un buen world model debe dar, antes del experimento, la dirección, el tamaño del efecto y la incertidumbre; cuando el resultado experimental no coincide, hay que registrarlo y actualizar el modelo. Mostrar solo mapas de calor exitosos y elegir a posteriori casos interpretables convierte un sistema científico en una demostración no falsable.
\end{importantbox}

\subsection{Por qué un simulador perfecto sigue necesitando un scientist o un agent}

Aunque el simulador pueda ejecutarse mil millones de veces a bajo costo, alguien tiene que definir qué intervenciones vale la pena simular. El espacio combinatorio es enorme, y muchos experimentos, aunque produzcan cambios, no aportan conocimiento nuevo o no pueden verificarse en la realidad. La selección de experimentos es un problema de active learning y de decision-making, no un problema de puro rendimiento de inferencia.

\teachervoice{01:06:01--01:07:34, el docente dice que, incluso con un simulador de pacientes perfecto, se necesita un agent que le indique qué vale la pena simular. Los expertos del dominio suelen tener ya un primer conjunto de hipótesis de alto valor; la función del simulador es compararlas y filtrarlas más rápido.}

\teachervoice{01:07:43--01:10:07, el docente mantiene una postura abierta pero mesurada hacia el AI scientist: la escala no es necesariamente el cuello de botella actual, y los agents existentes tampoco superan de forma significativa los experimentos que los expertos proponen con rapidez; aun así, vale la pena probar herramientas de ciclo cerrado.}

\begin{knowledgebox}{Límites razonables de las funciones de un AI scientist}
Los papeles más viables a corto plazo incluyen recuperar evidencia existente, proponer controles, invocar simuladores, cuantificar la incertidumbre, generar calendarios experimentales y registrar fracasos. Si puede proponer de forma independiente mecanismos de gran avance debe determinarse mediante evaluaciones ciegas, experimentos prospectivos y comparación con líneas base de expertos, no por la fluidez del texto que genera.
\end{knowledgebox}

Un ciclo cerrado con capacidad de auditoría debe incluir: versión de los datos, versión del modelo, intervenciones candidatas, intervalos de predicción, justificación de la selección, resultados experimentales y análisis de fracasos. Solo si los resultados negativos también se reincorporan al entrenamiento puede el sistema corregir gradualmente el world model, en lugar de amplificar continuamente el sesgo de confirmación.

\subsection{Términos clave}

Esta sesión abarca el aprendizaje automático y la biología; al final, una tabla vuelve a situar los términos de alta densidad en el problema que cada uno resuelve, para no quedarse solo con los nombres.

\begin{center}
\small
\begin{tabularx}{\textwidth}{p{0.18\textwidth} X X}
\toprule
\textbf{Término} & \textbf{Problema que resuelve} & \textbf{Mecanismo central y relación con la clase} \\
\midrule
World model & Cómo cambia el futuro después de una acción & Modela la distribución de estados futuros a partir de las observaciones multimodales actuales y de las acciones candidatas. \\
Translation & Cómo transferir los hechos compartidos entre modalidades & Usa una modalidad para predecir otra o alinea representaciones compartidas, como de H\&E a expresión génica. \\
Disambiguation & Cómo eliminar la ambigüedad de una sola modalidad & Usa la información exclusiva de otra modalidad para actualizar la posterior, como el vecindario espacial que ayuda a determinar el estado celular. \\
Cross-attention & Cómo interactúan de forma dirigida dos flujos de tokens & La query proviene del flujo objetivo y las key/value, del flujo de condición. \\
Adaptive LayerNorm & Cómo inyectar condiciones sin alargar la secuencia & Una red de condición genera la escala y el desplazamiento de la normalización y controla el cálculo de cada capa de la red troncal. \\
H\&E & Observar la morfología del tejido a bajo costo & Tinción clínica de rutina, de amplia cobertura pero con información molecular indirecta. \\
Spatial transcriptomics & Dónde ocurre la expresión génica & Conserva a la vez los niveles de expresión y las coordenadas del tejido, y conecta el estado molecular con la ecología espacial. \\
Imputation & Cómo estimar modalidades costosas o faltantes & Aprende la distribución condicional e informa la incertidumbre; no puede tomarse como sustituto sin pérdida de una medición real. \\
Counterfactual & Cómo cambia la predicción del modelo al modificar la entrada & Sirve para ordenar intervenciones candidatas; un cambio dentro del modelo no equivale a un efecto causal real. \\
Calibration & Si la confianza es creíble & Los intervalos de predicción deben tener la cobertura correcta en held-out patients y en distintos genes. \\
Information bottleneck & Cómo comprimir enormes cantidades de tokens multiescala & Elige representaciones a nivel de paciente, tejido, vecindario y célula, equilibrando eficiencia y conservación de señales poco frecuentes. \\
\bottomrule
\end{tabularx}
\end{center}

\subsection{Resumen de la sección}

Las preguntas y respuestas completan los límites reales del sistema: las muestras a nivel de paciente son limitadas, los controles sanos faltan de forma no aleatoria, los contrafactuales del modelo requieren validación experimental y el simulador no elige por sí solo las preguntas valiosas. El papel más razonable a corto plazo para un agent científico es organizar la evidencia y las herramientas de ciclo cerrado, no sustituir el diseño experimental con fluidez lingüística.

